\subsection{Liftings of homomorphisms in complete filtered algebras}

Let \(k\) be a ring, C a \(k\)-algebra, \((C_n)_{n \in \mathbf{Z}}\) a decreasing filtration of C, compatible with the structure of \(k\)-algebra and such that \(C_0 = C\) (III, § 2, no. 1). Suppose C is separated and complete for the topology defined by this filtration, so that the canonical map \(C \to \varprojlim C / C_n\) is a homeomorphism (loc. cit., no. 6). Let \(m\) be an integer \(>0\). Let us denote \(\pi: C \to C / C_m\) the canonical surjection.

PROPOSITION 5.—Let A be a formally smooth linearly topologized \(k\)-algebra. Every continuous homomorphism of \(k\)-algebras \(\varphi: A \to \mathrm{C/C}_{m}\) admits a continuous lifting to C.

For every integer \(n > m\), let us denote \(\pi_n: \mathrm{C/C}_n \to \mathrm{C/C}_{n-1}\) the canonical surjection. Since C is identified with the projective limit of the \(\mathrm{C/C}_n\), it is equivalent to give a continuous lifting of \(\varphi\) to C or a family \((\varphi_n)_{n > m}\) of continuous homomorphisms of \(k\)-algebras \(\varphi_n: \mathrm{A} \to \mathrm{C/C}_n\), satisfying \(\pi_n \circ \varphi_n = \varphi_{n-1}\). This brings us back, by induction on \(m\), to proving the statement when \(\mathrm{C}_{m+1} = 0\). The ideal \(\mathrm{C}_m\) is then of square zero (since \(2m \geqslant m+1\)), whence the proposition, since A is formally smooth.

Example.—Let C be a \(k\)-algebra and N a nilpotent ideal of C. The proposition applies to the algebra C equipped with the N-adic filtration. If A is a formally smooth linearly topologized \(k\)-algebra, we obtain that every continuous homomorphism from A into the \(k\)-algebra C/N, endowed with the discrete topology, lifts to a continuous homomorphism from A into the \(k\)-algebra C, equipped with the discrete topology.

